\chapter{Simulation Design and Results}
\label{chap:main-simulation}

The simulation compares the four procedures from
Chapter~\ref{chap:main-methods} on the same observations. Three models vary
either the slope or the predictor covariance, allowing us to examine how
these changes affect estimation and prediction. The comparison also separates
the behaviour of the raw and Arnoldi representations of FPLS. Their common
statistical target provides a useful way to assess numerical stability,
although each procedure still chooses its own dimension.

The design follows the principal settings of
\textcite{babii_carrasco_tsafack_2026}. The complete specification and extended
results are retained in Appendices~\ref{chap:simulation-design}
and~\ref{chap:estimation-results}; supplementary distributions appear in
Appendix~\ref{app:simulation-diagnostics}.

\section{Three controlled data-generating models}
\label{main:sec:simulation-design}

Curves are observed on $T=200$ equally spaced points
$s_t=(t-1)/(T-1)$ in $[0,1]$. All calculations use the endpoint-average inner
product $\langle f,g\rangle_T=T^{-1}\sum_{t=1}^T f(s_t)g(s_t)$.
There is no additional measurement error, smoothing, or interpolation.
The generating basis has $J=100$ functions,
\begin{equation}
    \phi_1(s)=1,\qquad
    \phi_j(s)=\sqrt{2}\cos(j\pi s),\quad j=2,\ldots,100.
    \label{main:eq:simulation-basis}
\end{equation}
The missing frequency $\cos(\pi s)$ is intentional: this is the basis used
in the experiment. Its functions are orthonormal in continuous $L^2[0,1]$,
but not exactly under the endpoint-average inner product. Consequently,
the quantities $\lambda_j$ below are the generating score variances, rather
than the exact eigenvalues of the finite-grid covariance operator.

Independent standard normal scores $Z_{ij}$ generate the observations as
\begin{equation}
    X_i(s_t)=\sum_{j=1}^{100}\sqrt{\lambda_j}Z_{ij}\phi_j(s_t),
    \qquad
    Y_i=\langle X_i,\beta\rangle_T+\varepsilon_i,
    \qquad \varepsilon_i\sim\mathcal N(0,1),
    \label{main:eq:simulation-dgp}
\end{equation}
where the errors and scores are mutually independent, and
$\beta(s)=\sum_{j=1}^{100}b_j\phi_j(s)$. Population means and the intercept
are zero. The observations are passed directly to the estimators without
subtracting their realised sample means. This convention is common to all
four methods.

Model~1 uses $b_j=4j^{-2.7}$ and $\lambda_j=2j^{-1.1}$. Model~2 retains
these score variances but sets $b_1=\cdots=b_5=4$, leaving the remaining
coefficients unchanged. Model~3 retains the slope of Model~1 but sets
$\lambda_1=\cdots=\lambda_5=2$, leaving the variance tail unchanged.
Since $b_1=4$ and $\lambda_1=2$ already hold in Model~1, the changes affect
only indices $2$--$5$. Model~2 therefore changes the target while preserving
the predictor distribution; Model~3 changes the predictor distribution while
preserving the target.

The response signal also changes. Its exact finite-grid variance is $\sum_j\lambda_j\langle\beta,\phi_j\rangle_T^2$, giving
$32.520$, $71.306$, and $33.061$ in Models~1--3. Because the error variance
is one, these are also the signal-to-noise ratios. Model~2 requires the
first five generating directions to account for $99\%$ of predictive
signal; Models~1 and~3 require only two. This describes the generating
model, not the number of components an estimated procedure should select.
In particular, raw MSPEs should primarily be compared between methods
within a model, because Model~2 has substantially more response variation.

Figure~\ref{main:fig:simulation-signal} shows how this signal is distributed.
Model~2 combines large slope coefficients in several directions with the
same declining score variances as Model~1. Its additional slope features
must therefore be estimated from directions with progressively less predictor
variation. Model~3 changes the information available about these directions
without introducing new slope features. The distinction matters because
neither the total size of the slope nor the predictor-variance ordering
alone determines which components are useful for predicting the response.
The complete slope and covariance plots are in
Appendix~\ref{chap:simulation-design}.

\begin{figure}[htbp]
    \centering
    \includegraphics[width=\textwidth]{predictive_signal_allocation.pdf}
    \caption[Predictive signal in the three simulation models]{Contribution
    $q_j=\lambda_j\langle\beta,\phi_j\rangle_T^2$ of each generating
    direction to response signal, and its cumulative share. Calculations
    use all $100$ directions under the endpoint-average inner product;
    the first twenty are displayed. Open markers identify the first
    component count exceeding $99\%$ of total signal.}
    \label{main:fig:simulation-signal}
\end{figure}

\section{Estimation, tuning, and evaluation}
\label{main:sec:simulation-evaluation}

Each model uses $5{,}000$ independent training samples of size $n=100$.
For every training sample, all methods receive the same data and are
evaluated on the same independent test sample of $100$ observations.
The three models have independent random streams, generated from the
fixed seed $2026$. Test observations play no part in tuning.

The candidate dimensions are $1,\ldots,70$. Raw and Arnoldi FPLS minimise
their own five-fold prediction cross-validation errors using identical
folds. Raw candidates must yield valid finite fits in the full sample
and every training fold. Arnoldi uses two-pass classical Gram--Schmidt
and a response-space QR solve, with rank tolerance $10^{-10}$. FPCR
minimises response-space GCV over eligible positive eigenvalues; the
smallest minimiser is selected when criteria tie. These choices compare
complete procedures, including their tuning, rather than estimates forced
to have equal dimension or effective degrees of freedom.

CG--FPLS uses the moment-residual objective and short recurrence specified
in Chapter~\ref{chap:main-methods}. It stops at the first component satisfying
\begin{equation}
    \|\widehat r-\widehat K\widehat\beta_m^{\mathrm{CG}}\|_T
    \leq 1.01\left(\frac{2\widehat\sigma^2\widehat E_X}{0.1n}\right)^{1/2},
    \qquad
    \widehat E_X=\frac1n\sum_{i=1}^{n}\|X_i\|_T^2.
    \label{main:eq:simulation-stopping}
\end{equation}
An internal FPCR fit chosen by moment-space GCV supplies
$\widehat\sigma^2$ through its response residual sum of squares divided
by $n$. This preliminary fit estimates the threshold scale; it is
different from the reported response-GCV FPCR procedure. If the threshold
is never reached, the rule returns $m=70$ and records this outcome.
The rule uses estimated quantities and is therefore not an oracle supplied
with the true noise variance.

Slope recovery and independent prediction are measured separately:
\begin{align}
    \operatorname{ISE}
    &=\|\widehat\beta-\beta\|_T^2,
    \label{main:eq:simulation-ise}\\
    \operatorname{MSPE}
    &=\frac1{100}\sum_{i=1}^{100}
      \{Y_i^{\mathrm{test}}-
      \langle X_i^{\mathrm{test}},\widehat\beta\rangle_T\}^{2}.
    \label{main:eq:simulation-mspe}
\end{align}
For a fixed fitted slope, expected MSPE equals
$1+\sum_j\lambda_j\langle\phi_j,\widehat\beta-\beta\rangle_T^2$.
The noise floor is therefore one. Prediction weights slope errors
according to predictor variation, whereas ISE measures the overall slope
error. This difference is central to interpreting the results.

Every finite error remains in the main analysis. A raw fit is flagged
when the condition number of its column-normalised Krylov basis or
response design reaches $\varepsilon_{\mathrm{mach}}^{-1/2}
\simeq6.71\times10^7$. Arnoldi is flagged if its orthogonality defect
reaches $\varepsilon_{\mathrm{mach}}^{1/2}\simeq1.49\times10^{-8}$.
Non-finite losses would also be flagged for every method. An error
exceeding $100$ times its own model--method positive median is labelled
extreme, but is never deleted for that reason. Means, medians, upper
percentiles, and maxima are reported together in the extended results.

\section{Slope recovery and prediction}
\label{main:sec:simulation-risk}

Table~\ref{main:tab:simulation-risk} contains all $60{,}000$ fitted results.
All losses are finite, including those associated with the $445$ flagged
raw-FPLS fits. The means include the impact of rare large errors; the
medians describe typical performance. For the three numerically stable
methods, Monte Carlo standard errors of mean ISE range from $0.0018$ to
$0.0475$, and those of mean MSPE from $0.0024$ to $0.0088$.

\begin{table}[htbp]
    \centering
    \small
    \caption[Main simulation estimation and prediction results]{Slope and prediction errors over all $5{,}000$ samples per
    model. Every finite error, including flagged raw-FPLS results, is retained.}
    \label{main:tab:simulation-risk}
    \begin{tabular}{@{}llrrrr@{}}
        \toprule
        & & \multicolumn{2}{c}{Slope ISE} & \multicolumn{2}{c}{Test MSPE}\\
        \cmidrule(lr){3-4}\cmidrule(l){5-6}
        Model & Method & Mean & Median & Mean & Median\\
        \midrule
        1 & CG--FPLS     & 0.633 & 0.597 & 1.270 & 1.230\\
          & Raw FPLS     & 334.298 & 0.539 & 21.224 & 1.114\\
          & Arnoldi FPLS & 0.760 & 0.539 & 1.130 & 1.113\\
          & FPCR         & 0.670 & 0.532 & 1.126 & 1.112\\
        \addlinespace
        2 & CG--FPLS     & 5.506 & 5.272 & 1.999 & 1.861\\
          & Raw FPLS     & 7.594 & 5.384 & 1.518 & 1.428\\
          & Arnoldi FPLS & 6.065 & 5.345 & 1.442 & 1.419\\
          & FPCR         & 5.662 & 5.365 & 1.495 & 1.472\\
        \addlinespace
        3 & CG--FPLS     & 0.447 & 0.426 & 1.221 & 1.176\\
          & Raw FPLS     & 1.219 & 0.446 & 1.154 & 1.114\\
          & Arnoldi FPLS & 0.670 & 0.446 & 1.131 & 1.114\\
          & FPCR         & 0.615 & 0.463 & 1.127 & 1.113\\
        \bottomrule
    \end{tabular}
\end{table}

CG--FPLS has the lowest mean ISE in all three models. This does not
mean that it always gives the most accurate typical estimate. In Model~1,
FPCR has the smallest median ISE, $0.532$, compared with $0.597$ for CG.
The CG distribution is more concentrated, however, giving a lower mean
ISE of $0.633$. Raw FPLS has a competitive median of $0.539$ but a mean
of $334.298$, showing how strongly an occasional numerical failure can
affect average performance.

The ordering differs for prediction. In Models~1 and~3, FPCR and Arnoldi
have mean MSPEs around $1.13$, whereas CG reaches $1.270$ and $1.221$.
In Model~2, where the slope is spread across more directions, Arnoldi
has the lowest mean MSPE, $1.442$, followed by FPCR at $1.495$.
CG has the lowest mean ISE, $5.506$, but the largest mean MSPE, $1.999$.
The reduction in total slope error therefore does not compensate for
errors in directions that matter more for prediction.

Model~2 is especially informative because the predictors are unchanged
from Model~1: the different results reflect the changed slope and the
procedures' responses to it. CG and Arnoldi have relatively close median
ISEs, $5.272$ and $5.345$, even though their median MSPEs are $1.861$
and $1.419$. The advantage in mean slope error is therefore insufficient
to identify the better predictor. Raw FPLS also has a competitive median
MSPE of $1.428$, but its mean rises to $1.518$ and its numerical flags
are more frequent. This setting brings together both issues: the loss
function changes the statistical comparison, while unreliable computation
can damage a procedure whose typical prediction is accurate.

Changing only the covariance in Model~3 improves CG mean ISE from
$0.633$ to $0.447$. More predictor variation provides more information
about the corresponding slope coefficients, while also increasing the
prediction penalty for errors in those directions. The two effects
explain why estimation and prediction need not improve by the same
amount. Across these models, no single ranking answers both questions.

\section{Selected dimensions and the noise estimate}
\label{main:sec:simulation-dimensions}

Table~\ref{main:tab:simulation-dimensions} shows the median selected
dimensions alongside raw-FPLS instability. The Krylov methods generally
use fewer components than FPCR. Their directions use response information,
whereas FPCR orders directions by predictor variation. One Krylov component
can combine several generating basis directions, so these counts cannot
be interpreted as a common measure of model complexity.

\begin{table}[htbp]
    \centering
    \small
    \caption[Selected dimensions and raw-FPLS stability]{Median selected dimensions and raw-FPLS instability among
    $5{,}000$ samples per model. No other method receives an instability flag.}
    \label{main:tab:simulation-dimensions}
    \begin{tabular}{@{}lrrrrrr@{}}
        \toprule
        & \multicolumn{4}{c}{Median dimension} & \multicolumn{2}{c}{Raw flags}\\
        \cmidrule(lr){2-5}\cmidrule(l){6-7}
        Model & CG & Raw & Arnoldi & FPCR & Count & Rate\\
        \midrule
        1 & 1 & 2 & 2 & 5  & 20  & $0.40\%$\\
        2 & 3 & 5 & 5 & 17 & 409 & $8.18\%$\\
        3 & 2 & 3 & 3 & 6  & 16  & $0.32\%$\\
        \bottomrule
    \end{tabular}
\end{table}

All $15{,}000$ CG fits satisfy the discrepancy rule before the upper
limit. CG selects one component in $90.2\%$ of Model~1 samples, three
in $88.6\%$ of Model~2 samples, and two in $83.2\%$ of Model~3 samples.
The internal noise estimate nevertheless has an important limitation.
Its preliminary moment-GCV fit selects a median of $70$ components in
every model and estimates a mean noise variance of approximately $0.300$,
well below the true value of one.

This is consistent with dividing an overfitted response residual sum
of squares by $n$ without a degrees-of-freedom adjustment:
$1-70/100=0.30$. It is not evidence of accurate noise estimation.
A smaller estimated noise level lowers the discrepancy threshold and,
along the same residual path, can only delay stopping. Thus the observed
one-to-three-component selections reflect a rapid residual reduction even
with this underestimated scale. The results describe the implemented
feasible procedure; they do not validate an oracle discrepancy rule.

The largest raw selections are $57$, $70$, and $68$ across Models~1--3,
compared with Arnoldi maxima of $14$, $16$, and $11$. A high selected
dimension alone cannot establish that additional signal has been recovered.
Its interpretation depends on the associated numerical diagnostics.

\section{Numerical stability}
\label{main:sec:simulation-stability}

Figure~\ref{main:fig:simulation-tails} shows the complete error
distributions, including the raw-FPLS extremes. Instability is most
frequent in Model~2, where more Krylov directions are usually selected.
Yet the less frequent failures in Model~1 produce the most severe average
loss: just one ISE exceeds $100$ times the raw median, reaching
$1.668\times10^6$. The largest raw MSPE is $1.005\times10^5$.
Its otherwise competitive median therefore gives an incomplete account
of reliability.

\begin{figure}[htbp]
    \centering
    \includegraphics[width=\textwidth]{babii_style_error_boxplots_with_tails.pdf}
    \caption[Full simulation error distributions]{Slope ISE (left) and independent test MSPE (right), with
    one row per model. All errors are included on logarithmic axes.
    Red open diamonds identify flagged raw-FPLS fits; whiskers use the
    conventional $1.5$ interquartile-range rule. CG, Raw, and Arnoldi
    denote the corresponding procedures in Chapter~\ref{chap:main-methods}.}
    \label{main:fig:simulation-tails}
\end{figure}

Arnoldi avoids these numerical failures in the experiment, but numerical
stability does not remove statistical variation. Its largest ISEs are
$36.35$, $77.41$, and $30.23$. These are not accompanied by loss of
basis orthogonality. The corresponding CG maxima are $1.82$, $16.79$,
and $1.25$, consistent with its stronger regularisation.

\section{Explaining the bias--variance trade-off}
\label{main:sec:simulation-bias-variance}

The average squared slope error separates exactly into squared bias and
variance across the simulated samples. Writing
$\overline\beta=R^{-1}\sum_{r=1}^{R}\widehat\beta_r$, with $R=5{,}000$,
\begin{equation}
    \frac1R\sum_{r=1}^{R}\|\widehat\beta_r-\beta\|_T^2
    =\|\overline\beta-\beta\|_T^2+
    \frac1R\sum_{r=1}^{R}\|\widehat\beta_r-\overline\beta\|_T^2.
    \label{main:eq:simulation-decomposition}
\end{equation}
For prediction, the corresponding bias and variance contributions are
weighted by $\lambda_j$ in the generating score directions. Adding the
noise variance of one gives expected prediction risk. This expectation
averages over a new observation; the observed test MSPE in
Table~\ref{main:tab:simulation-risk} also reflects the finite test samples,
so the two need not coincide exactly.

Table~\ref{main:tab:simulation-bias-variance} reports these components for
the three methods whose slope estimates can be reconstructed reliably
in every sample. All $5{,}000$ samples enter each row. Raw FPLS is omitted
from this decomposition because its $445$ flagged slopes are too sensitive
for reliable function-valued reconstruction; its full losses remain in
Table~\ref{main:tab:simulation-risk} and
Figure~\ref{main:fig:simulation-tails}.

\begin{table}[htbp]
    \centering
    \small
    \caption[Bias and variance across all simulation samples]{Squared bias
    and variance over all $5{,}000$ samples per model. Prediction columns
    exclude irreducible noise variance, which equals one. No samples are
    removed from the three displayed methods.}
    \label{main:tab:simulation-bias-variance}
    \begin{tabular}{@{}llrrrr@{}}
        \toprule
        & & \multicolumn{2}{c}{Slope risk}
          & \multicolumn{2}{c}{Excess prediction risk}\\
        \cmidrule(lr){3-4}\cmidrule(l){5-6}
        Model & Method & Sq.\ bias & Variance & Sq.\ bias & Variance\\
        \midrule
        1 & CG--FPLS     & 0.1069 & 0.5259 & 0.1102 & 0.1604\\
          & Arnoldi FPLS & 0.0082 & 0.7520 & 0.0132 & 0.1163\\
          & FPCR         & 0.0098 & 0.6602 & 0.0164 & 0.1102\\
        \addlinespace
        2 & CG--FPLS     & 1.2263 & 4.2796 & 0.4733 & 0.5224\\
          & Arnoldi FPLS & 0.0990 & 5.9656 & 0.0491 & 0.3903\\
          & FPCR         & 0.2419 & 5.4197 & 0.1106 & 0.3857\\
        \addlinespace
        3 & CG--FPLS     & 0.0183 & 0.4282 & 0.0346 & 0.1883\\
          & Arnoldi FPLS & 0.0055 & 0.6642 & 0.0098 & 0.1230\\
          & FPCR         & 0.0072 & 0.6078 & 0.0135 & 0.1141\\
        \bottomrule
    \end{tabular}
\end{table}

The decomposition helps explain CG's advantage in mean ISE.
In Model~1, CG slope variance
is $0.526$, compared with $0.752$ for Arnoldi, while squared bias is
$0.107$ rather than $0.008$. The variance reduction outweighs the extra
bias in the slope norm. Under prediction weighting, however, CG has both
larger squared bias ($0.110$ versus $0.013$) and larger variance
($0.160$ versus $0.116$). Model~2 shows the same distinction more strongly:
CG's prediction squared bias is $0.473$, compared with $0.049$ for
Arnoldi. Early stopping therefore need not reduce variance under every
loss measure.

In Model~2, CG reduces slope variance relative to Arnoldi by approximately
$1.69$, more than compensating for the increase of approximately $1.13$
in squared bias. Prediction weights change that balance: both CG's
prediction bias and its prediction variance exceed the Arnoldi values.
The smaller number of CG components should therefore not be interpreted
as a uniform reduction in uncertainty. The benefit depends on which
directions are uncertain and how the loss function weights them.

An additional comparison retains only the $4{,}980$, $4{,}591$, and
$4{,}984$ samples where the raw fit passes its stability check. All four
methods use these same samples for the paired comparison. This restriction
supports reliable reconstruction of slopes and their bias--variance
components; it never replaces the full ISE or MSPE results.
On this subset, raw and Arnoldi mean MSPEs are almost identical:
$1.1279$ versus $1.1282$, $1.4271$ versus $1.4282$, and $1.1302$
for both in Model~3. Each method still chooses its own dimension,
so this is not a comparison at a fixed common $m$.

These findings agree with their exact-arithmetic equivalence at a
common valid dimension. Reorthogonalisation and QR make the same
response-least-squares target more reliably computable as the raw powers
become nearly collinear. They do not guarantee smaller statistical error
on every sample.

The coefficient distributions in Appendix~\ref{app:simulation-diagnostics}
provide a more direct view of the slope differences. In Model~1, CG
estimates the second coefficient with mean $0.366$, compared with its
true value $0.616$, and the third with mean $0.071$, compared with
$0.206$. In Model~3, greater predictor variation brings these means to
$0.599$ and $0.198$. The same slope features are estimated more accurately
when the predictors contain more information about them. Raw-FPLS
distribution panels use only its stable fits, whereas the other methods
use all samples. These supplementary plots illustrate particular
coefficients and locations; their normal quantile plots do not establish
Gaussianity of the full estimated function.

The conclusions apply to these three models, dense error-free curves,
Gaussian scores, $n=100$, and the stated tuning rules. They establish
neither a universal ranking nor equal-complexity comparisons.
Chapter~\ref{chap:main-empirical} examines observed spectra.
Chapter~\ref{chap:main-inference} first considers a different requirement:
for the stated testing argument, the remaining CG residual must be small
relative to sampling uncertainty. A dimension suitable for estimation
need not satisfy that requirement.
